\chapter{Results and Discussion}
\label{chap:results_and_discussion}

\hspace*{\parindent}This chapter presents the results of the development, testing, and evaluation of KABISIG. In accordance with the testing and evaluation procedures established in Chapter~4 (Section~4.2.18.2), the findings are organized across system implementation outcomes, functional and security verification, cross-browser and performance measurements, ISO/IEC~25010:2023 technical product quality assessment, Technology Acceptance Model (TAM) user acceptance results, qualitative stakeholder feedback, and a synthesis of conclusions, recommendations, and limitations.

\section{System Implementation and Deployment Summary}
\label{sec:system_implementation_and_deployment_summary}

\hspace*{\parindent}The KABISIG web application was developed using an Agile/Scrum iterative life cycle spanning seven months and deployed in a cloud-hosted production environment on Vercel (frontend and serverless API hosting), backed by Supabase (PostgreSQL) for relational storage, authentication, and tenant-scoped file storage. All five core modules specified in the project objectives were implemented and operational during field testing:

\begin{enumerate}
    \item \textbf{Module 1 (Kabataan Profiling and Digital Registry):} public youth registration with age validation (15--30 years), SK Chairperson ID verification workflows, automated Digital Youth ID generation with unique QR codes, deterministic demographic classification, and rule-based Youth Engagement Scoring.
   \newpage
    \item \textbf{Module 2 (Program, Event, and Rule-Based Recommendation Management):} program CRUD operations, integrated calendar views, schedule and venue conflict detection, verification-gated registration, and browser camera-based QR attendance check-in (\texttt{html5-qrcode} and \texttt{getUserMedia}).
    \item \textbf{Module 3 (Financial Transparency and Governance Support):} annual budget setup aligned with DILG JMC No.~1, s.~2025, itemized expense recording with automated VAT/non-VAT computations, asset inventory tracking, governance document approval routing, compliance deadline monitoring, and the read-only Public Transparency Dashboard.
    \item \textbf{Module 4 (Community Engagement, Communication, and Rule-Based Feedback Analytics):} the \textit{Boses ng Kabataan} portal with optional anonymity (\texttt{user\_id} set to \texttt{NULL}), keyword-dictionary sentiment classification, program request extraction, automated spam and profanity flagging, community polls, resolution voting, and in-app notifications.
    \item \textbf{Module 5 (Analytics, Reporting, and Public Information):} barangay operational dashboards, Super Admin cross-barangay federation analytics across all 27 pre-provisioned barangay tenants in Naga City, client-side COA-ready PDF report generation (\texttt{jsPDF}), and official Facebook Page publishing via Meta Graph API v25.0.
\end{enumerate}
\vspace{5mm}

\section{Respondent Profile and Turnout}
\label{sec:respondent_profile_and_turnout}

\hspace*{\parindent}The sampling plan (Section~4.2.18.3) projected 73 to 75 respondents. Actual evaluation data collection yielded 35 respondents (about 47--48\% of the target) across four user groups representing the core operational roles of the platform:

\begin{itemize}
    \item \textbf{SK Officials ($n = 15$):} 6 Chairpersons, 2 Kagawads, 4 Secretaries, and 3 Treasurers from participating barangays in Naga City, who evaluated administrative workflows during live testing sessions, including the 36th SK Federation Regular Session on November~16, 2026, and the final system testing at the Naga City SK Session Hall.
    \item \textbf{Youth Constituents ($n = 17$):} registered youth aged 15 to 30 who tested public portal registration, Digital Youth ID viewing, program registration, polling, and \textit{Boses ng Kabataan} feedback submission.
    \item \textbf{Public Viewers ($n = 2$):} unauthenticated community members who evaluated the login-free Public Transparency Dashboard, financial summaries, and public announcements.
    \item \textbf{IT Expert / Technical Evaluator ($n = 1$):} an independent IT practitioner with 13 years of experience, not involved in system development, who conducted a technical quality audit using the ISO/IEC~25010:2023 instrument.
\end{itemize}

\hspace*{\parindent}Actual turnout fell short of the 73 to 75 target due to overlapping Sangguniang Kabataan council schedules, limited youth availability during the evaluation window, and tight pre-defense academic deadlines. SK officials reached half of their target (15 of 30), youth reached 57\% of theirs (17 of 30), public viewers reached 20\% of the 10-respondent target, and only one of the three to five planned IT evaluators participated. Equal quota stratification across youth subgroups was not fully achieved. All 35 respondents agreed to the informed consent statement, which follows the Data Privacy Act of 2012 (RA~10173). Results are reported per respondent group using descriptive statistics (mean).
\vspace{5mm}

\begin{table}[H]
\centering
\small
\caption{Distribution of Evaluation Respondents ($N = 35$)}
\label{tab:distribution_evaluation_respondents}
\begin{tabular}{|p{3.5cm}|c|c|c|p{3.6cm}|}
\hline
\multicolumn{1}{|c|}{\textbf{Respondent Group}} & \textbf{Target} & \textbf{Actual ($n$)} & \textbf{Percentage (\%)} & \multicolumn{1}{c|}{\textbf{Instrument Administered}} \\ \hline
SK Officials (Chairperson, Kagawad, Secretary, Treasurer) & 30 & 15 & 42.86 & UAT + Role-Based TAM Questionnaire \\ \hline
Youth Constituents (Ages 15--30) & 30 & 17 & 48.57 & UAT + Youth TAM Questionnaire \\ \hline
Public Viewers (Unauthenticated Citizens) & 10 & 2 & 5.71 & Public Portal UAT + Public Viewer TAM \\ \hline
IT Experts / Technical Evaluators & 3--5 & 1 & 2.86 & ISO/IEC~25010:2023 Quality Instrument \\ \hline
\textbf{Total} & \textbf{73--75} & \textbf{35} & \textbf{100.00} & \textbf{Mixed Quantitative and Qualitative Survey} \\ \hline
\end{tabular}
\end{table}
\vspace{5mm}

\begin{table}[H]
\centering
\small
\caption{Respondent Characteristics and Modules Tested ($N = 35$)}
\label{tab:respondent_characteristics_modules_tested}
\begin{tabular}{|p{4.0cm}|p{4.2cm}|c|c|}
\hline
\multicolumn{1}{|c|}{\textbf{Characteristic}} & \multicolumn{1}{c|}{\textbf{Category}} & \textbf{$f$} & \textbf{\%} \\ \hline
Age range & 15 to 20 & 4 & 11.4 \\ \cline{2-4}
 & 21 to 25 & 24 & 68.6 \\ \cline{2-4}
 & 26 to 30 & 6 & 17.1 \\ \cline{2-4}
 & N/A (IT expert) & 1 & 2.9 \\ \hline
Prior experience with digital or online systems (multiple answers allowed) & None & 1 & 2.9 \\ \cline{2-4}
 & Basic & 15 & 42.9 \\ \cline{2-4}
 & Moderate & 19 & 54.3 \\ \cline{2-4}
 & Extensive & 2 & 5.7 \\ \hline
Modules tested (multiple answers allowed) & Module 1: Profiling and Digital Registry & 22 & 62.9 \\ \cline{2-4}
 & Module 2: Programs, Events, and Recommendations & 8 & 22.9 \\ \cline{2-4}
 & Module 3: Financial Transparency and Governance & 14 & 40.0 \\ \cline{2-4}
 & Module 4: Community Engagement and Feedback & 16 & 45.7 \\ \cline{2-4}
 & Module 5: Analytics, Reporting, and Public Information & 31 & 88.6 \\ \hline
\end{tabular}
\end{table}
\vspace{2mm}

\noindent\textit{Note.} Percentages are based on $N = 35$. The prior-experience and module rows allow multiple answers, so their frequencies exceed the number of respondents and their percentages exceed 100\%.
\vspace{3mm}

\hspace*{\parindent}Most respondents were aged 21 to 25 and reported basic to moderate digital experience (see Figures~\ref{fig:appa_respondent_group}--\ref{fig:appa_modules_tested}). Evaluation coverage was uneven across modules. Module 5 was tested by nearly nine in ten respondents, while Module 2 was tested by fewer than a quarter. Findings about program and QR attendance management therefore rest on a smaller group of evaluators than findings about dashboards and reporting.
\vspace{5mm}

\section{Functional, Security, Performance, and Compatibility Testing Results}
\label{sec:functional_security_performance_and_compatibility_testing_results}

\subsection{Functional and Security Testing}
\label{subsec:functional_and_security_testing}

\hspace*{\parindent}Prior to stakeholder UAT, all five modules underwent structured functional and security test cases verifying CRUD operations, Role-Based Access Control (RBAC), PostgreSQL Row-Level Security (RLS) multi-tenant isolation by \texttt{tenant\_id}, age and registration gating, duplicate QR attendance prevention, financial tax calculations, and audit log recording.
\vspace{5mm}

\begin{table}[H]
\centering
\footnotesize
\setlength{\tabcolsep}{4pt}
\caption{Summary of Functional and Security Test Cases}
\label{tab:summary_functional_and_security_test_cases}
\begin{tabular}{|p{2.9cm}|c|c|c|c|p{5.5cm}|}
\hline
\multicolumn{1}{|c|}{\textbf{Module / Domain}} & \textbf{Executed} & \textbf{Passed} & \textbf{Failed} & \textbf{Pass (\%)} & \multicolumn{1}{c|}{\textbf{Remarks}} \\ \hline
Module 1: Kabataan Profiling and Digital ID & 8 & 8 & 0 & 100.00\% & Public registration, age validation (15--30), barangay residency check, ID upload, Chairperson approval/rejection, Digital Youth ID \& QR generation, rule-based demographic classification, and engagement scoring. \\ \hline
Module 2: Programs and QR Attendance & 8 & 8 & 0 & 100.00\% & Program CRUD, calendar integration, schedule/venue conflict detection, verification-gated registration, slot capacity checks, camera QR check-in, and duplicate scan rejection. \\ \hline
Module 3: Financial Transparency and Documents & 9 & 9 & 0 & 100.00\% & Annual budget setup, program fund allocation, over-budget warning alerts, itemized expense entry, automated VAT/non-VAT \& remaining balance math, inventory CRUD, document upload/approval routing, compliance status tracking, and public transparency filtering. \\ \hline
Module 4: Boses ng Kabataan and Polls & 7 & 7 & 0 & 100.00\% & Identified feedback submission, anonymous toggle (\texttt{user\_id = NULL}), keyword sentiment tagging, program request extraction, spam/profanity auto-flagging, community poll voting, and resolution voting. \\ \hline
Module 5: Analytics, PDF Export, and FB Publishing & 6 & 6 & 0 & 100.00\% & Barangay operational charts, Super Admin cross-barangay aggregation, client-side COA PDF generation (\texttt{jsPDF}), in-app notifications, and Meta Graph API Facebook Page posting. \\ \hline
RBAC, Multi-Tenant RLS, and Audit Logging & 7 & 7 & 0 & 100.00\% & Role-based route restrictions, cross-tenant \texttt{tenant\_id} query blocking via PostgreSQL RLS, Zod payload validation, signed URL document access, leadership transfer demotion/promotion, and \texttt{audit\_logs} recording. \\ \hline
\textbf{Overall} & \textbf{45} & \textbf{45} & \textbf{0} & \textbf{100.00\%} & \multicolumn{1}{c|}{---} \\ \hline
\end{tabular}
\end{table}
\vspace{5mm}

\subsection{Performance and Cross-Browser Compatibility Testing}
\label{subsec:performance_and_cross_browser_compatibility_testing}

\hspace*{\parindent}Timing benchmarks were recorded during deployment testing against the thresholds defined in Section~4.2.18.2 (page load time under 2.0~s and REST API latency under 500~ms), alongside supplementary benchmark targets for client-side COA PDF generation and QR camera scan recognition. Notification dispatch speed, which Section~4.2.18.2 also lists, was not separately benchmarked.
\vspace{5mm}

\begin{table}[H]
\centering
\footnotesize
\setlength{\tabcolsep}{4pt}
\caption{System Performance Benchmarks}
\label{tab:system_performance_benchmarks}
\begin{tabular}{|>{\raggedright\arraybackslash}p{0.20\linewidth}|>{\raggedright\arraybackslash}p{0.20\linewidth}|>{\centering\arraybackslash}p{0.11\linewidth}|>{\centering\arraybackslash}p{0.09\linewidth}|>{\raggedright\arraybackslash}p{0.29\linewidth}|}
\hline
\textbf{Performance Metric} & \textbf{Target Threshold} & \textbf{Measured Average} & \textbf{Status} & \textbf{Supporting Evidence (ISO Evaluation and Feedback)} \\ \hline
Initial Dashboard / Page Load Time & $< 2.00\text{ s}$ (Sec.~4.2.18.2) & $1.15\text{ s}$ & Passed & Page loads were rated Good by the IT evaluator under Performance Efficiency. \\ \hline
REST API Response Latency (Express + Supabase) & $< 500\text{ ms}$ (Sec.~4.2.18.2) & $240\text{ ms}$ & Passed & REST API responses were rated Good by the IT evaluator. \\ \hline
Client-Side COA PDF Generation (\texttt{jsPDF}) & $< 2.00\text{ s}$ (supplementary target, not defined in Sec.~4.2.18.2) & $1.45\text{ s}$ & Passed & Client-side generation avoids server round-trips; part of an item rated Fair overall by the IT evaluator alongside QR scanning and concurrent multi-tenant queries. \\ \hline
QR Camera Scan Recognition (\texttt{html5-qrcode}) & $< 1.50\text{ s}$ (supplementary target, not defined in Sec.~4.2.18.2) & $1.65\text{ s}$ & Not met & Rated Fair by the IT evaluator, with respondents noting slower scanning/camera switching in low-light or weak-signal venues. \\ \hline
Notification Dispatch Speed & No preset target (Sec.~4.2.18.2) & Not measured & --- & Listed in Section~4.2.18.2; not separately benchmarked during testing. \\ \hline
\end{tabular}
\end{table}
\vspace{5mm}

\hspace*{\parindent}Cross-browser and responsive layout verification covered desktop browsers (Google Chrome, Mozilla Firefox, Microsoft Edge, and Apple Safari) and mobile browsers (Android Chrome and iOS Safari) using TailwindCSS utility layouts, without requiring a native mobile application. In the ISO/IEC 25010:2023 evaluation (Section~\ref{sec:iso_25010_software_product_quality_evaluation}), the independent evaluator rated cross-browser operation across desktop and mobile browsers as Good, whereas responsive layout rendering across desktop, tablet, and smartphone screens received a Fair rating (contributing to a Compatibility characteristic mean of 2.50). Accordingly, responsive layout adaptation across varying screen sizes and viewports warrants further refinement and testing prior to wider rollout.
\vspace{5mm}

\section{ISO/IEC~25010:2023 Software Product Quality Evaluation}
\label{sec:iso_25010_software_product_quality_evaluation}

\hspace*{\parindent}Technical product quality was assessed by the independent IT evaluator ($n = 1$) across eight product quality characteristics of ISO/IEC~25010:2023 (excluding Safety), using the 4-point scale in Appendix~F ($4 = \text{Very Good}$, $3 = \text{Good}$, $2 = \text{Fair}$, $1 = \text{Poor}$; see Figures~\ref{fig:appa_functional_suitability}--\ref{fig:appa_portability}). The 2023 edition renamed Usability as Interaction Capability and Portability as Flexibility, and these names are used below. Each characteristic mean is the average of its item ratings.
\vspace{5mm}

\begin{table}[H]
\centering
\footnotesize
\setlength{\tabcolsep}{4pt}
\caption{ISO/IEC~25010:2023 Technical Evaluation Summary ($n = 1$)}
\label{tab:iso_25010_technical_evaluation}
\begin{tabular}{|p{2.8cm}|c|c|c|p{4.2cm}|}
\hline
\multicolumn{1}{|c|}{\textbf{Quality Characteristic}} & \textbf{Ratings} & \textbf{Mean} & \textbf{Interpretation} & \multicolumn{1}{c|}{\textbf{Key Technical Observations}} \\ \hline
1. Functional Suitability & 3, 3, 3 & 3.00 & Good & All five modules, budget and VAT/non-VAT computations, and rule-based features were rated Good. \\ \hline
2. Performance Efficiency & 3, 2 & 2.50 & Fair & Page loads and REST responses were rated Good. PDF generation, QR scanning, and concurrent multi-tenant queries were rated Fair. \\ \hline
3. Compatibility & 3, 2 & 2.50 & Fair & Operation across desktop and mobile browsers was rated Good. Responsive layouts across desktop, tablet, and smartphone screens were rated Fair. \\ \hline
4. Interaction Capability (Usability) & 3, 2, 3 & 2.67 & Good & Role-based menus and validation prompts were rated Good. Dashboards, charts (\texttt{Recharts}/\texttt{D3.js}), and data tables were rated Fair. \\ \hline
5. Reliability & 3, 4 & 3.50 & Very Good & Runtime stability was rated Good. Data integrity across sessions and \texttt{Supabase Storage} file uploads were rated Very Good. \\ \hline
6. Security & 4, 4, 4 & 4.00 & Very Good & RBAC and \texttt{tenant\_id} RLS isolation, protection of administrative actions, and protection of youth profiles and anonymous feedback were all rated Very Good. \\ \hline
7. Maintainability & 4, 3 & 3.50 & Very Good & The modular three-tier architecture was rated Very Good. Updating tenants, committees, and roles without disrupting live records was rated Good. \\ \hline
8. Flexibility (Portability) & 4, 4 & 4.00 & Very Good & Zero-install browser access and cloud hosting across operating systems were rated Very Good. \\ \hline
\textbf{Overall Mean} & --- & \textbf{3.21} & \textbf{Good} & Approaches the Very Good threshold of 3.26 \\ \hline
\end{tabular}
\end{table}
\vspace{3mm}

\noindent\textbf{Scale:} 3.26--4.00 = Very Good; 2.51--3.25 = Good; 1.76--2.50 = Fair; 1.00--1.75 = Poor.
\vspace{5mm}

\hspace*{\parindent}The overall mean of 3.21 falls in the Good range, just below the Very Good threshold. Security and Flexibility received the highest rating (4.00), followed by Reliability and Maintainability (3.50 each). These results indicate that the multi-tenant Row-Level Security architecture, \texttt{Zod 4.5.4} schema validation, and the separation across presentation, application, and data tiers provide a strong foundation for multi-barangay deployment.
\vspace{5mm}

\hspace*{\parindent}Performance Efficiency and Compatibility each averaged 2.50, which is the upper limit of the Fair range. The lower ratings concern client-side PDF generation (\texttt{jsPDF}), camera QR scanning (\texttt{html5-qrcode}), concurrent multi-tenant use, responsive layouts, and dashboard presentation. These ratings are consistent with the respondents' comments in Section~5.6 asking for offline caching of QR scans and faster scanning in weak-signal venues. Although client-side PDF generation met its supplementary 2.00~s target, the evaluator's Fair rating applied to a combined item covering PDF generation, QR scanning, and concurrent multi-tenant use; QR scan recognition was the only measured metric that missed its target. Because the ratings come from a single evaluator, they describe one expert's judgment and should not be generalized. Further audits by several independent evaluators are recommended.
\vspace{5mm}

\section{Technology Acceptance Model (TAM) Evaluation Results}
\label{sec:technology_acceptance_model_evaluation_results}

\hspace*{\parindent}End-user acceptance was measured across SK Officials ($n = 15$), Youth Constituents ($n = 17$), and Public Viewers ($n = 2$), for 34 end users in total (see Figures~\ref{fig:appa_chairperson_pu}--\ref{fig:appa_public_viewer_biu}). Role-tailored TAM instruments (Appendix~F, Part~2) used a 5-point Likert scale ($5 = \text{Strongly Agree}$ to $1 = \text{Strongly Disagree}$) across four constructs: Perceived Usefulness (PU), Perceived Ease of Use (PEOU), Attitude Toward Use (ATU), and Behavioral Intention to Use (BIU). Group means were computed from all item responses within each construct. The overall column weights each group by its number of respondents, and the composite is the average of the four construct means.
\vspace{5mm}

\begin{table}[H]
\centering
\footnotesize
\setlength{\tabcolsep}{4pt}
\caption{Summary of TAM Evaluation Results by Respondent Group ($n = 34$ End Users)}
\label{tab:summary_tam_evaluation_results_by_respondent_group}
\begin{tabular}{|p{3.2cm}|c|c|c|c|c|}
\hline
\multicolumn{1}{|c|}{\textbf{TAM Construct}} & \textbf{SK Officials} & \textbf{Youth} & \textbf{Public} & \textbf{Overall} & \textbf{Interpretation} \\
 & \textbf{($n = 15$)} & \textbf{($n = 17$)} & \textbf{($n = 2$)} & \textbf{Mean} & \\ \hline
Perceived Usefulness (PU) & 4.28 & 3.87 & 4.63 & 4.10 & Agree \\ \hline
Perceived Ease of Use (PEOU) & 4.32 & 3.76 & 4.63 & 4.06 & Agree \\ \hline
Attitude Toward Use (ATU) & 4.33 & 3.55 & 4.50 & 3.95 & Agree \\ \hline
Behavioral Intention to Use (BIU) & 4.69 & 3.53 & 4.83 & 4.12 & Agree \\ \hline
\textbf{Composite TAM Mean} & \textbf{4.41} & \textbf{3.68} & \textbf{4.65} & \textbf{4.06} & \textbf{Agree} \\ \hline
\textbf{Group Interpretation} & \textbf{Strongly Agree} & \textbf{Agree} & \textbf{Strongly Agree} & --- & --- \\ \hline
\end{tabular}
\end{table}
\vspace{3mm}

\noindent\textbf{Scale:} 4.21--5.00 = Strongly Agree; 3.41--4.20 = Agree; 2.61--3.40 = Neutral; 1.81--2.60 = Disagree; 1.00--1.80 = Strongly Disagree.
\vspace{5mm}

\begin{table}[H]
\centering
\small
\caption{TAM Results of SK Officials by Role}
\label{tab:tam_results_of_sk_officials_by_role}
\begin{tabular}{|l|c|c|c|c|c|c|}
\hline
\multicolumn{1}{|c|}{\textbf{Role}} & \textbf{$n$} & \textbf{PU} & \textbf{PEOU} & \textbf{ATU} & \textbf{BIU} & \textbf{Composite} \\ \hline
Chairperson & 6 & 4.21 & 4.42 & 4.39 & 4.56 & 4.39 \\ \hline
Kagawad & 2 & 4.50 & 5.00 & 4.50 & 5.00 & 4.75 \\ \hline
Secretary & 4 & 4.19 & 4.06 & 4.17 & 4.50 & 4.23 \\ \hline
Treasurer & 3 & 4.42 & 4.00 & 4.33 & 5.00 & 4.44 \\ \hline
\textbf{All SK Officials} & \textbf{15} & \textbf{4.28} & \textbf{4.32} & \textbf{4.33} & \textbf{4.69} & \textbf{4.41} \\ \hline
\end{tabular}

\vspace{2mm}
\par\noindent Note. Interpretations are based on unrounded means.
\end{table}
\vspace{3mm}

\hspace*{\parindent}Across all three groups, no respondent chose Disagree or Strongly Disagree on any TAM item. Overall acceptance was positive (Agree, 4.06), but it differed by group.
\vspace{5mm}

\subsection{Discussion of TAM Findings by Stakeholder Group}
\label{subsec:discussion_of_tam_findings_by_stakeholder_group}

\begin{itemize}
    \item \textbf{SK Officials ($n = 15$):} Officials rated all four constructs in the Strongly Agree range, with the highest score for Behavioral Intention to Use (4.69). They noted that automated VAT and non-VAT expense calculations, QR attendance check-ins (\texttt{html5-qrcode} and \texttt{getUserMedia}), centralized document approval routing, and instant COA-ready PDF exports (\texttt{jsPDF}) remove repetitive manual spreadsheet and paper-based work. By role, Kagawads and Treasurers gave the strongest intention-to-use ratings, while Treasurers gave the lowest ease-of-use ratings (4.00), with responses evenly split across Strongly Agree, Agree, and Neutral on financial workflows and error frequency. Barangay Administrators (Chairpersons) gave a few Neutral answers on annual budget setup (aligned with DILG JMC No.~1, s.~2025) and the barangay dashboard.
    \item \textbf{Youth Constituents ($n = 17$):} Youth rated every construct in the Agree range (composite 3.68). They valued online registration, mobile access to their Digital Youth ID QR code, and the option to send anonymous suggestions through \textit{Boses ng Kabataan}. Neutral responses were common, especially for Attitude Toward Use and Behavioral Intention. For ``I enjoy using KABISIG'' (3.35) and ``I would recommend KABISIG'' (3.47), about two thirds or more of youth chose Neutral. Possible explanations include short testing sessions, limited exposure to some modules, and the absence of an immediate personal need for the features. These explanations were not tested in this study, and follow-up interviews with youth are recommended.
    \item \textbf{Public Viewers ($n = 2$):} Public viewers recorded high ratings across all four TAM constructs---Perceived Usefulness (4.63), Perceived Ease of Use (4.63), Attitude Toward Use (4.50), and Behavioral Intention to Use (4.83), with a composite score of 4.65 (Strongly Agree)---reflecting favorable impressions of login-free access to public financial and program information, though with only two respondents these results are descriptive and cannot be generalized.
\end{itemize}
\vspace{5mm}

\section{Qualitative Feedback Synthesis}
\label{sec:qualitative_feedback_synthesis}

\hspace*{\parindent}Open-ended responses in Part~3 of the instrument (Appendix~F; see Figures~\ref{fig:appa_qual_strengths_1}--\ref{fig:appa_qual_other_2}) drew 10 responses on strengths, 7 on needed improvements, 6 on errors or confusing steps, and 12 on other comments. Several answers were ``None'' or ``N/A''. Comments written in Filipino and Bikol were translated and grouped by theme by the researchers.
\vspace{5mm}

\subsection{Identified System Strengths}
\label{subsec:identified_system_strengths}

\begin{itemize}
    \item \textbf{Centralization:} an ``all-in-one'' platform that replaces scattered chats and folders for documents, resolutions, and youth profiles, and covers the day-to-day workflow of SK councils.
    \item \textbf{Faster, less manual work:} QR attendance tracking (\texttt{html5-qrcode} and \texttt{getUserMedia}) was described as much faster than manual paper lists, and respondents reported a large reduction in manual encoding for program creation and attendance.
    \item \textbf{Financial automation and transparency:} automatic VAT/non-VAT and remaining-budget computation (avoiding formula mistakes in Excel), easier liquidation, and a public view of how the barangay budget is allocated and spent via the Public Transparency Dashboard.
    \item \textbf{Youth voice and information access:} \textit{Boses ng Kabataan} helps officials stay aware of youth comments, respondents praised the feedback and poll summaries for quickly showing which projects appeal to youth, and the system acts as a bridge between youth and their barangay while providing reliable information on barangay projects and services.
    \item \textbf{Technical design:} the IT evaluator observed that KABISIG goes beyond basic CRUD by embedding domain rules in its workflows, and that the separation of duties among Barangay Administrator (Chairperson), Kagawad, Secretary, and Treasurer mirrors real local governance controls.
    \item \textbf{Adoption interest:} one respondent noted that their barangay already uses its own system but that KABISIG has more features and may be adopted, and others hoped it would be pushed to the council as their working system.
\end{itemize}
\vspace{5mm}

\subsection{Suggestions and Areas for Improvement}
\label{subsec:suggestions_and_areas_for_improvement}

\begin{itemize}
    \item \textbf{Offline and low-signal support:} lightweight offline caching or a local sync queue for QR scanning, and faster camera switching and scanning in dim or weak-signal venues such as covered courts.
    \item \textbf{Audit and approval tools:} quick filters for the audit log and approval queue as records pile up, and an append-only \texttt{audit\_logs} table at database level so historical records stay tamper-proof during leadership transitions.
    \item \textbf{Announcements:} a designated person to post and update upcoming barangay programs and announcements.
    \item \textbf{Youth data fields:} one respondent suggested that the system stop requiring year-level and scholarship status inputs for graduates.
    \item \textbf{Errors reported:} most responses reported no errors, aside from one partial comment mentioning the user's QR ID, which should be checked against the QR issuance logs.
\end{itemize}
\vspace{5mm}

\section{Summary of Findings, Conclusions, and Recommendations}
\label{sec:summary_of_findings_conclusions_and_recommendations}

\subsection{Summary of Findings}
\label{subsec:summary_of_findings}

\begin{enumerate}
    \item \textbf{Constraints of Existing SK Operations (Objective 1):} Prior to KABISIG, SK councils in Naga City relied on disconnected spreadsheets, physical folders, and social media chats (Section~4.2.16.1 and Appendix~B), resulting in delayed reporting, manual tax computation errors, duplicate beneficiary tracking issues, and limited public financial visibility.
    \item \textbf{Multi-Tenant Platform Development (Objectives 2 and 3):} KABISIG was engineered as a cloud-hosted, three-tier multi-tenant web system using Next.js, React, TypeScript, TailwindCSS, Express 5.2.1 (Node.js), and Supabase (PostgreSQL with Row-Level Security), integrating all five governance modules provisioned for Naga City's 27 barangays.
    \item \textbf{User Acceptance and Technical Quality (Objective 4):} Across 35 evaluation respondents (15 SK officials, 17 youth constituents, 2 public viewers, and 1 IT expert), KABISIG achieved a composite TAM rating of 4.06 (Agree) among the 34 end users. SK officials rated it 4.41 (Strongly Agree), youth 3.68 (Agree), and public viewers 4.65 (Strongly Agree). The IT expert's ISO/IEC~25010:2023 rating was 3.21 (Good), with Security and Flexibility rated highest (4.00) and Performance Efficiency and Compatibility lowest (2.50).
    \item \textbf{System Performance and Reliability (Objective 5):} Functional and security testing yielded a 100.00\% pass rate (45 of 45 test cases passed) across the core modules and RBAC/RLS isolation checks. Measured performance averages were 1.15~s for initial dashboard/page load time (meeting $< 2.00\text{ s}$), 240~ms for REST API response latency (meeting $< 500\text{ ms}$), 1.45~s for client-side COA PDF generation (\texttt{jsPDF}; within the $< 2.00\text{ s}$ supplementary target), and 1.65~s for QR camera scan recognition (\texttt{html5-qrcode}), which did not meet the $< 1.50\text{ s}$ supplementary target. The independent evaluator rated PDF generation, QR scanning, and concurrent-use responsiveness as Fair, and respondents requested offline support for QR attendance.
\end{enumerate}
\vspace{5mm}

\subsection{Conclusions}
\label{subsec:conclusions}

\hspace*{\parindent}Based on the development outcomes and evaluation results, the researchers conclude that:
\vspace{3mm}

\begin{enumerate}
    \item KABISIG can modernize and standardize barangay-level youth governance in Naga City by replacing fragmented manual workflows with a unified, role-based digital platform.
    \item Multi-tenant data isolation by \texttt{tenant\_id} and role-based access controls support secure barangay-level operations and federation-level oversight within the pre-provisioned schema, as reflected in the independent evaluator's Very Good ratings across Security items.
    \item User acceptance varied across stakeholder groups---highest among public viewers and SK officials, and more moderate among youth---indicating that role-aligned operational workflows and public transparency features are well received while youth engagement requires further follow-up.
    \item Technical quality is Good overall. Reliability, Security, Maintainability, and Flexibility are strengths, while Performance Efficiency, responsive-layout compatibility, and dashboard presentation need improvement before city-wide deployment.
    \item All 45 functional and security test cases passed, and page load, API latency, and PDF generation met their targets, while QR scan recognition (1.65~s) did not meet its supplementary target and concurrent-load performance was not tested.
\end{enumerate}
\vspace{5mm}

\subsection{Recommendations}
\label{subsec:recommendations}

\hspace*{\parindent}Aligned with the study's scope limitations (Section~1.4.2), actual respondent turnout, and the future enhancement roadmap (Section~4.2.18.6), the following recommendations are proposed:
\vspace{3mm}

\begin{enumerate}
    \item \textbf{Multi-Barangay Institutional Adoption:} coordinate formal turnover and onboarding for the 22 active Sangguniang Kabataan councils with the Super Administrator (SK Federation President) and the Naga City ICT Office, supported by periodic training for newly elected or appointed officers.
    \item \textbf{Offline Capability and Scanner Improvements:} implement Progressive Web App (PWA) offline caching or a local sync queue for QR attendance check-ins, alongside faster camera switching and recognition in low-light or weak-signal venues, addressing respondent feedback and the Fair rating on Performance Efficiency.
    \item \textbf{Audit Log and Approval Queue Enhancements:} add quick filters to the audit log and approval queue for high-volume records, and enforce append-only behavior on the \texttt{audit\_logs} table at the database level to preserve historical integrity during leadership transitions.
    \item \textbf{Youth Engagement and Data Quality:} investigate why youth responses were predominantly Neutral (such as through follow-up interviews), make year-level and scholarship or graduation status fields optional for graduates, and designate a clear owner for posting announcements.
    \item \textbf{Responsive Layout and Dashboard Refinement:} conduct targeted testing and refinement of responsive layouts across desktop, tablet, and smartphone viewports, as well as the visual clarity of dashboard charts and tables.
    \item \textbf{Expanded Longitudinal Evaluation:} carry out a follow-up multi-barangay evaluation after one full fiscal cycle, including a load test across multiple tenants, with larger youth and public viewer samples and multiple independent IT evaluators to assess long-term governance impact and statistical generalizability.
    \item \textbf{Roadmap Enhancements:} evaluate the future enhancements in Section~4.2.18.6, including SMS and email notification gateways, hard spending limits with mandatory justification text for over-budget allocations and expenses, automatic linking of itemized expenses to inventory entries, a scholar tracking module, integration with Naga City government services and national government databases (such as DILG, PSA, and NYC), and predictive analytics for budget forecasting and program outcome evaluation.
\end{enumerate}
\vspace{5mm}

\subsection{Limitations}
\label{subsec:limitations}

\hspace*{\parindent}The evaluation reached 35 of the 73 to 75 planned respondents, with one IT evaluator and two public viewers, so ratings for those groups are descriptive only. Module 2 was tested by only 8 respondents. Performance was measured under normal use, and no load or concurrency test was run, so scalability across all 27 provisioned tenants is not yet verified.
\vspace{5mm}